\section{Supplement to Section 3: authority condition}\label{supp:feas}

This section gives the settings and secondary results of Section~\ref{M-sec:feasibility}.

\subsection{Reproduction of the published condition}\label{supp:feas-pub}

The condition of Eq.~\eqref{M-eq:pub} was evaluated with the simulator of \ref{M-app:simulator} without the clip, with single regeneration and over 50 revolutions, starting with a step entry from rest at the depth $\alim(g)$; the demand is $K_ag\max_t\lvert\dot x\rvert$ over the whole run, compared with $F_{\max}=\SI{27}{\newton}$. The gain was searched on $\mathcal G=[0,3000]$\,\si{\volt\second\per\metre} by two methods: the self-adaptive differential evolution (SADE) of \citet{qin2009sade}, with 20 individuals, at most 60 generations, a stop after 25 generations without improvement and three seeds per speed; and a deterministic reference, a scan in steps of \SI{20}{\volt\second\per\metre} refined by bisection at the edges of every feasible interval. Both evaluate the same objective, and the map takes the better result.
% src: src/objective.py (demand: K_a g max|x'|, single regeneration, N_REVS = 50), src/sade.py, src/gainmap.py (G = [0, 3000], DG = 20, reference), src/run_gain_map.py; README §2, §4.6

SADE reached the optimum at 250 of the 251 speeds. Where it and the reference differ (2280, 2300 and \SI{2310}{\rpm}), one of them missed an isolated pocket of feasible gains: the \SI{20}{\volt\second\per\metre} grid of the reference stepped over the pocket at 2300 and \SI{2310}{\rpm}, and all three seeds of SADE missed it at \SI{2280}{\rpm} (\SI{5.96}{\milli\metre} against \SI{12.03}{\milli\metre}), a property of the objective rather than a numerical error. The optimum lies in such a pocket, above an interval of rejected gains, at every speed from 2060 to \SI{2310}{\rpm} and at 670, 1010 and 1970--\SI{2000}{\rpm} (32 speeds). On the interval connected to $g=0$ the gain at 2060--\SI{2310}{\rpm} would be about 0.72 times the published one, and the depth 6.0--\SI{8.0}{\milli\metre}. Where the optimum enters or leaves a pocket the map jumps, for instance from 383 to \SI{213}{\volt\second\per\metre} between 2000 and \SI{2010}{\rpm} and from 858 to \SI{415}{\volt\second\per\metre} between 2310 and \SI{2320}{\rpm}. Over all speeds the reproduced gain ranges from \SI{71}{\volt\second\per\metre} (\SI{1950}{\rpm}) to \SI{1046}{\volt\second\per\metre}.
% src: README §5.2 (250/251; pockets; connected-interval gain 0.72 of the published); results/gainmap_main.npz (g_ref, a_ref, g_sade, a_sade, sade_runs: 2280 rpm all seeds 5.96 mm, a_ref 12.03 mm; 2300 and 2310 rpm a_ref 5.97 and 5.99 mm, SADE 11.87 and 11.80 mm; scan_g, scan_f, scan_alim, scan_fd: isolated pockets recomputed, 32 speeds; connected-interval depths recomputed with the edge bisected as in gainmap.reference: 5.96-7.96 mm, against 5.75-7.52 mm at the last feasible gain of the 20 V s/m grid; g_best 71.1-1045.7 V s/m, minimum at 1950 rpm, 383.3 -> 213.0 and 858.5 -> 414.8 V s/m)

The ratios to the published map quoted in Section~\ref{M-sec:feas-pub} use the map of Fig.~8 of \citet{ozsoy2026robust}, read from the vector graphics of that figure. The published map is straight over 500--550, 560--610 and 1410--\SI{1860}{\rpm} and is a single segment without vertices over 620--\SI{740}{\rpm}; these segments are excluded, so that the medians 0.74 and 0.88 refer to 750--1150 and 1160--\SI{1400}{\rpm}. The median ratio at 2060--\SI{2310}{\rpm} is 1.83. If the demand $g\lvert\dot x\rvert$ is read as a voltage and compared with 27, the reproduced gains become 2--3 times larger (780--\SI{1971}{\volt\second\per\metre}).
% src: verification/targets/fig8_gain_map.csv, fig8_segments.csv; README §4.4 (straight segments), §5.2, §6 (voltage reading: results/gainmap_volt.npz, g_ref 780-1971 V s/m); results/paper_vs_ours.json (p04: band medians 0.74, 0.88, 1.83, 0.97; band_note: 500-1150 is 750-1150 and 1160-1490 is 1160-1400)

With the step entry the peak demand is that of the entry transient (Section~\ref{M-sec:feas-pub}). On a grid of \SI{50}{\volt\second\per\metre} the optimum of Eq.~\eqref{M-eq:pub} at \SI{2500}{\rpm} is \SI{500}{\volt\second\per\metre} ($\alim=\SI{5.86}{\milli\metre}$) with the step entry, but \SI{2900}{\volt\second\per\metre} (\SI{38.3}{\milli\metre}) with a feed ramp of one revolution, \SI{3000}{\volt\second\per\metre} (\SI{40.0}{\milli\metre}) with a feed ramp of ten revolutions and \SI{2950}{\volt\second\per\metre} (\SI{39.2}{\milli\metre}) with a face entry.
% src: verification/paper/support_runs.json (eq2_optimum_2500rpm: step 500 V s/m, 5.86 mm; ramp1 2900, 38.29 mm; ramp10 3000, 40.03 mm; face 2950, 39.15 mm)

\subsection{The reformulated condition}\label{supp:feas-new}

By Eqs.~\eqref{M-eq:Gcl} and~\eqref{M-eq:xp}, the harmonic of $b(t)$ at $\omega=m\omega_T$ is
\begin{equation}\label{eq:supp-feas-bm}
  b_m=-g\,i\omega\,G_{\mathrm{cl}}(i\omega)\,F_m(a)=-F_m(a)\,i\omega G(i\omega)\,\frac{g}{1+gz},
\end{equation}
with $z=i\omega G_a(i\omega)G(i\omega)$. At a fixed depth, $\lvert g/(1+gz)\rvert^{-2}=(1/g+\operatorname{Re}z)^2+(\operatorname{Im}z)^2$ has the derivative $-2(1/g+\operatorname{Re}z)/g^2$ with respect to $g$, so $\lvert b_m\rvert$ in Eq.~\eqref{eq:supp-feas-bm} grows with $g$ if and only if $1+g\operatorname{Re}z>0$: at every gain for the harmonics with $\operatorname{Re}z\ge0$, and for $g<1/\lvert\operatorname{Re}z\rvert$ at the others. The steady demand $F_{\mathrm{ss}}$ of Eq.~\eqref{M-eq:feas} can therefore fall with $g$ only where this condition fails, through the phases of the harmonics, which also change with $g$ and move the peak of their sum, or through the depth range $(0,\kappa\alim(g)]$ and the 80 depths sampled evenly in it, which move with $g$. On a grid of 151 gains from 0 to \SI{3000}{\volt\second\per\metre} at all 251 speeds, $F_{\mathrm{ss}}$ increased with $g$ except for drops of at most \SI{0.22}{\newton} at 27 speeds, all at $g\ge\SI{2120}{\volt\second\per\metre}$ and at demands above \SI{140}{\newton}, far beyond $F_{\max}$, and the feasible set was one interval $[0,g_b]$ at every speed. Its upper edge was found by bisection to \SI{0.25}{\volt\second\per\metre}; SADE, applied to the same condition, improved on it by at most \SI{0.013}{\milli\metre}.
% src: Eqs. (Gcl), (xp) of the article; src/structure.py (G, Ga, Gcl), src/feasibility.py (forced_velocity, demand_envelope: N_DEPTH = 80; edge_gain: tol 0.25); verification/paper/feas_monotone.json (251 speeds x 151 gains: drops <= 0.215 N at 27 speeds, g >= 2120 V s/m, F_ss >= 144.9 N; all_one_interval true); README §8.2 (SADE better by <= 0.013 mm); results/feas_g1500.npz (a_sade - a_ref <= 0.0133 mm)

The closed form was compared with the settled demand of the simulator at the operating depth $\kappa\alim$, at the 139 speeds where that run settled. At the 129 of them whose demand is at least \SI{1}{\newton}, the simulator, with 30 axial layers, exceeds the closed form by $-0.75$ to $+\SI{8.7}{\percent}$ (median $+\SI{1.1}{\percent}$; more than \SI{3.7}{\percent} at 12 speeds), by at most \SI{0.98}{\newton} (\SI{0.33}{\volt}); at the other ten, where the helix cancels the force harmonics, the demand is below \SI{1}{\newton} and the difference at most \SI{0.06}{\newton}. At three speeds (520, 660 and \SI{1900}{\rpm}, including the two largest differences), 240 layers bring the difference within \SI{0.8}{\percent}. The difference is therefore the axial quadrature error of the simulator, whose layers are thicker than the axial wavelengths of the higher force harmonics, whereas the closed form integrates over the depth exactly.
% src: results/feas_g1500.npz (cert_stable, cert_demand_td, cert_demand_cf: 129/139 points with demand >= 1 N, -0.75 % to +8.66 %, median +1.13 %, 12 above 3.7 %, at most 0.98 N; the other 10: demand < 0.92 N, |difference| <= 0.062 N); verification/feasibility/cf_vs_td_layers.json (520, 660, 1900 rpm: +8.66, +6.97, +3.73 % with 30 layers; -0.03, -0.13, -0.79 % with 240); README §8.2

The entry check of Eq.~\eqref{M-eq:achk} simulates step entries from rest with the clipped command ($\lvert V\rvert\le\SI{9}{\volt}$) and the full surface memory, each judged over 300 revolutions by the classifier D1 (\ref{M-app:simulator}). The depth starts at $0.225\alim=0.25\kappa\alim$ and grows by factors of 1.08 up to $\kappa\alim$; the first interval that ends in chatter is bisected to a relative \SI{0.3}{\percent}. With $\mathcal G=[0,1500]$ no shallower depth chattered where the operating depth settled, so the check binds exactly at the 112 speeds where the run at $\kappa\alim$ did not settle: at 85 speeds with $g^\star=\SI{1500}{\volt\second\per\metre}$ and at 1530--1560, 1980--2010 and 2280--\SI{2460}{\rpm}. At \SI{2500}{\rpm}, for instance, $\alim=\SI{17.24}{\milli\metre}$ and $\kappa\alim=\SI{15.51}{\milli\metre}$, but the clipped step entry chatters above \SI{13.01}{\milli\metre}. Over all speeds $\achk/\alim$ ranges from 0.74 to 0.90.
% src: src/run_feas_certify.py (a_lo = 0.25 kappa a_lim, growth 1.08, nreg = 0), src/tdboundary.py (boundary: rtol 3e-3; classify: NREV 300); results/feas_cert_g1500.npz, results/feas_g1500.npz (cert_stable = 0 at 112 speeds = 85 at g = 1500 + 27 at 1530-1560, 1980-2010, 2280-2460 rpm; no speed with cert_stable = 1 and a_cert < a_op_linear; 2500 rpm: a_lim 17.24, a_op_linear 15.51, a_td_clipped 13.01 mm; a_cert/a_lim 0.737-0.900); README §8.3

\subsection{Results by speed and sensitivity}\label{supp:feas-res}

In the range 1500--\SI{2050}{\rpm}, where the uncontrolled limit is already high and the forced resonance caps the gain, the mean $\achk$ (\SI{20.51}{\milli\metre}) is below the mean depth of Eq.~\eqref{M-eq:pub} (\SI{21.24}{\milli\metre}), although $\alim$ at the reformulated gains averages \SI{22.94}{\milli\metre} there; speed by speed, $\achk$ lies below the depth of Eq.~\eqref{M-eq:pub} at 49 of the 52 speeds from 1690 to \SI{2200}{\rpm}. Table~\ref{tab:supp-feas-table3} compares the two designs on an equal footing at the nine speeds of Table~3 of \citet{ozsoy2026robust}, with the same clipped step entry, the full surface memory and the same classifier.
% src: results/feas_cert_g1500.npz (a_cert, a_lim), results/gainmap_main.npz (a_best); README §8.4; per-speed count recomputed from these files: 49 of 52 at 1690-2200 rpm

\begin{table}[!tbp]
  \centering
  \caption{Clipped step-entry limit $a_{\mathrm{TD}}$ (\si{\milli\metre}; \SI{9}{\volt} clip, full surface memory) at the nine speeds of Table~3 of \citet{ozsoy2026robust}, with the gains $g$ (\si{\volt\second\per\metre}) of the published condition as reproduced [Eq.~\eqref{M-eq:pub}] and of the reformulated condition [Eq.~\eqref{M-eq:feas}, $\mathcal G=[0,1500]$].}
  \label{tab:supp-feas-table3}
  \begin{tabular}{@{}rrrrr@{}}
    \toprule
    & \multicolumn{2}{c}{Eq.~\eqref{M-eq:pub}} & \multicolumn{2}{c}{Eq.~\eqref{M-eq:feas}}\\
    \cmidrule(lr){2-3}\cmidrule(l){4-5}
    Speed (rpm) & $g$ & $a_{\mathrm{TD}}$ & $g$ & $a_{\mathrm{TD}}$\\
    \midrule
    500  & 829  & 12.1 & 1500 & 17.0\\
    560  & 1043 & 10.9 & 1500 & 16.2\\
    600  & 990  & 11.1 & 1500 & 15.7\\
    800  & 1026 & 10.9 & 1500 & 15.5\\
    1300 & 992  & 11.1 & 1500 & 14.3\\
    2300 & 849  & 12.1 & 1085 & 13.5\\
    2500 & 536  & 6.37 & 1500 & 13.0\\
    2800 & 668  & 7.33 & 1500 & 12.6\\
    3000 & 712  & 8.12 & 1500 & 12.8\\
    \bottomrule
  \end{tabular}
\end{table}
% src: results/table3_main.csv (g_prop, a_prop_td), results/feas_compare_g1500.json (table3: g_new, a_new_td); both with the 9 V clip and the full surface memory (feasibility.NREG = 0)

The margin $\kappa$ hardly matters: with $\kappa=0.95$ the gains change by at most \SI{0.53}{\volt\second\per\metre}, and the same 157 speeds sit at the bound. A guard of $0.85F_{\max}$ (\SI{23}{\newton}) lowers the gain at 116 speeds, 22 of which sat at the bound with a demand of 23.1--\SI{26.8}{\newton}; the gain then ranges from 101 to \SI{1500}{\volt\second\per\metre}, with 135 speeds at the bound. The cap of $\mathcal G$ matters more: with $g_{\max}=\SI{3000}{\volt\second\per\metre}$ the steady demand binds at 195 speeds and the gain reaches the bound at 56; the mean checked depth rises to \SI{17.1}{\milli\metre}, and the entry check binds at 184 speeds. At seven of these (570--\SI{990}{\rpm}) a shallower depth, 0.55--0.69 of the operating depth, chatters while the operating depth settles, which is why a ladder of depths is simulated.
% src: results/feas_k095.npz (at_bound 157, the same speeds as feas_g1500.npz; gains within 0.53 V s/m), results/feas_guard085.npz (116 speeds lowered, 22 of them at the bound in feas_g1500.npz with demand_ref 23.1-26.8 N; g_best 100.7-1500 V s/m; at_bound 135), results/feas_g3000.npz (at_bound 56; demand binds at 195), results/feas_cert_g3000.npz (mean a_cert 17.13 mm; a_cert < a_op_linear at 184 speeds; at 570, 580, 590, 630, 770, 780 and 990 rpm the run at a_op_linear settled, cert_stable = 1, while a_td_clipped/a_op_linear = 0.55-0.69); README §8.5, §8.6
